\chapter{Elliptic Functions on \texorpdfstring{\(\C/\Lambda\)}{C/Lambda}}
\label{ch:elliptic-functions}

A meromorphic function on \(T=\C/\Lambda\) is the same thing as a meromorphic
function \(f\) on \(\C\) satisfying \(f(z+\lambda)=f(z)\) for every
\(\lambda\in\Lambda\).  Such functions are called elliptic functions.

\begin{definition}
An elliptic function for \(\Lambda\) is a meromorphic function
\(f:\C\to\CP^1\) that is \(\Lambda\)-periodic.
\end{definition}

Compactness immediately rules out holomorphic examples.

\begin{proposition}
\label{prop:no-holomorphic-elliptic}
A holomorphic elliptic function is constant.
\end{proposition}

\begin{proof}
It descends to a holomorphic function on the compact surface \(T\), so
\cref{prop:compact-functions} applies.
\end{proof}

The useful elliptic functions must therefore have poles.  The exam residue
method now gives three constraints.  In one period parallelogram, an elliptic
function has residue sum zero, as many zeros as poles counted with multiplicity,
and the weighted sum of its zeros minus the weighted sum of its poles lies in
\(\Lambda\).

\begin{theorem}[Elliptic Liouville laws]
\label{thm:elliptic-liouville}
Let \(f\) be a nonzero elliptic function.  In a fundamental parallelogram avoiding
zeros and poles on its boundary,
\[
  \sum \Res(f\,dz)=0,\qquad
  \sum \ord_a(f)=0,\qquad
  \sum_{\ord_a(f)>0}\ord_a(f)a-\sum_{\ord_b(f)<0}(-\ord_b(f))b\in\Lambda .
\]
\end{theorem}

\begin{proof}[Proof sketch]
Integrate \(f\,dz\), \(f'/f\,dz\), and \(z f'/f\,dz\) around the boundary of a
period parallelogram.  Periodicity makes opposite sides cancel in the first two
integrals.  In the third integral, opposite sides cancel up to periods, which is
why the final answer is only well-defined modulo \(\Lambda\).
\end{proof}

\begin{keyidea}{The bridge to Abel}
The third law is already Abel's theorem in its torus form: a principal divisor of
degree zero has Abel sum zero in \(\C/\Lambda\).
\end{keyidea}

\section*{Exercises}
\addcontentsline{toc}{section}{Exercises}

\exerciseintent{Use the residue obstruction to rule out a single simple pole on a torus.}
\begin{exercise}
Assume a nonconstant elliptic function has total pole order \(1\).  Write the
Laurent expansion at its unique simple pole, compute the residue contribution,
and use the residue-sum law to obtain a contradiction.  Conclude that the total
pole order is at least \(2\).
\end{exercise}

\exerciseintent{Compute boundary cancellation on a period parallelogram via the argument principle.}
\begin{exercise}
Apply the argument principle on a fundamental parallelogram \(P\):
\[
  \frac{1}{2\pi i}\int_{\partial P}\frac{f'}{f}\,dz.
\]
Pair the four side integrals using periodicity and compute why the total is
\(0\).  Conclude that the number of zeros equals the number of poles, counted
with multiplicity.
\end{exercise}

\exerciseintent{See why Abel sums naturally live modulo the lattice.}
\begin{exercise}
Let a lift \(a_i\in\C\) of a point of \(T\) be replaced by \(a_i+\lambda_i\),
with \(\lambda_i\in\Lambda\).  Compute how
\(\sum_i m_i a_i\) changes, and conclude that the Abel sum is naturally a point
of \(\C/\Lambda\), not of \(\C\).
\end{exercise}

\exerciseintent{Verify the descent/ascent between elliptic functions and meromorphic functions on the quotient.}
\begin{exercise}
Let \(q:\C\to T=\C/\Lambda\) be the quotient map.  Show that a meromorphic
function \(F\) on \(T\) pulls back to an elliptic function by computing
\((F\circ q)(z+\lambda)\).  Conversely, if \(f(z+\lambda)=f(z)\), compute why
the formula \(F([z])=f(z)\) is independent of the chosen representative.
\end{exercise}

\exerciseintent{Compute an explicit principal-divisor Abel sum and connect it to \(\smap(\divisor(f))=0\).}
\begin{exercise}
Let \(f\) be a nonzero elliptic function with divisor
\[
  \divisor(f)=[a_1]+[a_2]-[b_1]-[b_2]
\]
on \(T\), with chosen lifts \(a_1,a_2,b_1,b_2\in\C\).  Use the third identity
in \cref{thm:elliptic-liouville} to compute
\[
  a_1+a_2-b_1-b_2\pmod{\Lambda}.
\]
Then rewrite the calculation as \(\smap(\divisor(f))=0\).
\end{exercise}
